% =====================================================================
%  Cours simple : LE TRI RAPIDE ET LE TRI PAR FUSION
%  CPGE deuxieme annee --- filieres MP / PSI / TSI
%  Compilation :  latexmk -pdf cours-rapide-fusion.tex
% =====================================================================
\documentclass[11pt,aspectratio=169]{beamer}
\usepackage{theme-cpge}

\begin{document}

% ---------------------------------------------------------------- titre
\begin{frame}[plain]
  \vfill
  \begin{center}
    {\color{teal}\hrule height 2.5pt}\vspace{12pt}
    {\Huge\bfseries\color{teal} Le tri rapide et le tri par fusion}\\[12pt]
    {\large Diviser pour régner}\\[6pt]
    {\normalsize Tri en place \;$\bullet$\; Stabilité \;$\bullet$\; Complexité}\\[14pt]
    {\normalsize CPGE deuxième année --- filières MP / PSI / TSI}\\[10pt]
    {\color{teal}\hrule height 1pt}
  \end{center}
  \vfill
\end{frame}

% ---------------------------------------------------------------- programme
\begin{frame}{Objectifs}
\begin{encadre}
\small
Cette section fait suite à celle étudiée en première année. L'objectif est d'explorer
des algorithmes de tri plus efficaces, tel que le \motcle{tri rapide} et le
\motcle{tri par fusion}, et de mettre l'accent sur la notion de
\motcle{tri en place} et de \motcle{stabilité}.
\end{encadre}
\vspace{8pt}
\begin{center}\small
\begin{tabular}{p{3cm}p{8cm}}
\toprule
\textbf{Tri rapide} & Définir le principe de l'algorithme. L'implémenter.
Comparer sa complexité avec celles des autres algorithmes. Discuter le choix du pivot.\\
\midrule
\textbf{Tri par fusion} & Définir le principe de la méthode. L'implémenter.
Comparer sa complexité avec celles des autres méthodes.\\
\bottomrule
\end{tabular}
\end{center}
\end{frame}

% ---------------------------------------------------------------- rappels
\begin{frame}{Trois mots de vocabulaire}
\begin{encadre}
\textbf{Tri en place} --- l'algorithme n'utilise qu'une mémoire auxiliaire
$\Oh(1)$ ou $\Oh(\log n)$ : il réorganise le tableau sans en fabriquer une copie.
\end{encadre}
\vspace{4pt}
\begin{encadre}
\textbf{Tri stable} --- deux éléments de même clé restent dans leur ordre d'origine.
Indispensable dès qu'on trie successivement sur plusieurs critères.
\end{encadre}
\vspace{4pt}
\begin{encadre}
\textbf{Complexité} --- on compte le nombre de \emph{comparaisons} entre éléments,
dans le meilleur cas, dans le pire cas, et en moyenne.
\end{encadre}
\vspace{6pt}
\centering\small
Rappel de première année : sélection, insertion et bulles coûtent
$\Th(n^2)$ comparaisons dans le pire cas.
\end{frame}

% ---------------------------------------------------------------- DPR
\begin{frame}{L'idée commune : diviser pour régner}
\begin{encadre}
\begin{itemize}\setlength\itemsep{2pt}
  \item \motcle{Diviser} : découper le tableau en deux parties ;
  \item \motcle{Régner} : trier chaque partie \emph{récursivement} ;
  \item \motcle{Combiner} : reconstituer le tableau trié.
\end{itemize}
\end{encadre}
\vspace{6pt}
\begin{center}\small
\begin{tabular}{lll}
\toprule
& \textbf{Diviser} & \textbf{Combiner}\\
\midrule
Tri par fusion & facile : couper au milieu & difficile : fusionner ($\Th(n)$)\\
Tri rapide & difficile : partitionner ($\Th(n)$) & rien à faire\\
\bottomrule
\end{tabular}
\end{center}
\vspace{6pt}
\centering\motcle{Les deux algorithmes se partagent le travail de façon opposée.}
\end{frame}

% =====================================================================
\section{Le tri par fusion}
% =====================================================================

\begin{frame}{Principe}
\begin{encadre}
Fusionner deux listes \emph{déjà triées} est facile : on compare les deux premiers
éléments et on prend le plus petit. D'où l'algorithme :
\textbf{couper en deux, trier chaque moitié, fusionner}.
\end{encadre}
\vspace{4pt}
\begin{center}
\begin{tikzpicture}[scale=0.8,transform shape,font=\scriptsize,
  b/.style={draw=teal,fill=tealclair,inner sep=3pt,rounded corners=1.5pt},
  r/.style={draw=vertrec,fill=vertrec!12,inner sep=3pt,rounded corners=1.5pt},
  >=Stealth]
  \node[b] (a0) at (0,2.2) {6 \ 2 \ 8 \ 4};
  \node[b] (a1) at (-2,1.2) {6 \ 2}; \node[b] (a2) at (2,1.2) {8 \ 4};
  \node[b] (b1) at (-3,0.2) {6}; \node[b] (b2) at (-1,0.2) {2};
  \node[b] (b3) at (1,0.2) {8}; \node[b] (b4) at (3,0.2) {4};
  \node[r] (c1) at (-2,-0.9) {2 \ 6}; \node[r] (c2) at (2,-0.9) {4 \ 8};
  \node[r] (d0) at (0,-1.9) {2 \ 4 \ 6 \ 8};
  \draw[->,teal] (a0)--(a1); \draw[->,teal] (a0)--(a2);
  \draw[->,teal] (a1)--(b1); \draw[->,teal] (a1)--(b2);
  \draw[->,teal] (a2)--(b3); \draw[->,teal] (a2)--(b4);
  \draw[->,vertrec] (b1)--(c1); \draw[->,vertrec] (b2)--(c1);
  \draw[->,vertrec] (b3)--(c2); \draw[->,vertrec] (b4)--(c2);
  \draw[->,vertrec] (c1)--(d0); \draw[->,vertrec] (c2)--(d0);
  \node[teal] at (-4.6,1.7) {\textbf{diviser}};
  \node[vertrec] at (-4.6,-1.4) {\textbf{fusionner}};
\end{tikzpicture}
\end{center}
\end{frame}

\begin{frame}[fragile]{L'étape de fusion}
\begin{lstlisting}
def fusion(u, v):
    """u et v sont triees ; renvoie la liste triee de tous leurs elements."""
    res = []
    i, j = 0, 0
    while i < len(u) and j < len(v):
        if u[i] <= v[j]:
            res.append(u[i]); i += 1
        else:
            res.append(v[j]); j += 1
    return res + u[i:] + v[j:]
\end{lstlisting}
\begin{encadre}\small
Chaque tour de boucle place définitivement un élément dans \tab{res} :
la fusion de deux listes totalisant $n$ éléments coûte au plus $n-1$ comparaisons,
donc $\Th(n)$.
\end{encadre}
\end{frame}

\begin{frame}[fragile]{L'algorithme}
\begin{lstlisting}
def tri_fusion(t):
    n = len(t)
    if n <= 1:                      # cas de base
        return t
    m = n // 2                      # diviser
    return fusion(tri_fusion(t[:m]), # regner
                  tri_fusion(t[m:])) # puis combiner
\end{lstlisting}
\vspace{2pt}
\begin{terminaisonrec}
La taille $|t|$ est un entier positif qui décroît strictement à chaque appel
($m<n$ et $n-m<n$) : on atteint nécessairement le cas de base $n\le1$.
\end{terminaisonrec}
\begin{correction}[Correction (récurrence sur $n$)]
\small
Les deux moitiés sont triées par hypothèse de récurrence, et la fusion de deux listes
triées renvoie une liste triée contenant exactement les mêmes éléments.
\end{correction}
\end{frame}

\begin{frame}{Complexité}
Soit $C(n)$ le nombre de comparaisons. La division coûte $0$, la fusion coûte au plus $n$ :
\[
C(n)\;=\;2\,C\!\left(\frac n2\right)+\Th(n),\qquad C(1)=0 .
\]
\vspace{-6pt}
\begin{center}
\begin{tikzpicture}[font=\scriptsize,yscale=0.9]
  \foreach \k/\lab/\c in {0/{1 appel de taille $n$}/{$n$},
                          1/{2 appels de taille $n/2$}/{$n$},
                          2/{4 appels de taille $n/4$}/{$n$}}{
    \draw[teal,fill=tealclair] (0,-\k*0.95) rectangle (6.4,-\k*0.95-0.55);
    \node at (3.2,-\k*0.95-0.28) {\lab};
    \node[anchor=west,teal] at (6.8,-\k*0.95-0.28) {coût : \c};
  }
  \node at (3.2,-2.9) {$\vdots$};
  \draw[teal,fill=tealclair] (0,-3.6) rectangle (6.4,-4.15);
  \node at (3.2,-3.88) {$n$ appels de taille $1$};
  \node[anchor=west,teal] at (6.8,-3.88) {coût : $n$};
  \draw[<->,thick,orangefonce] (-0.4,-0.28)--(-0.4,-3.88);
  \node[orangefonce,anchor=east] at (-0.6,-2) {$\log_2 n$};
\end{tikzpicture}
\end{center}
\vspace{-2pt}
\begin{encadre}
\centering
$\log_2 n$ niveaux, chacun de coût $\Th(n)$ \;$\Longrightarrow$\;
$C(n)=\Th(n\log n)$, \textbf{dans tous les cas}.
\end{encadre}
\end{frame}

\begin{frame}{Tri par fusion : bilan}
\begin{center}\small
\begin{tabular}{ll}
\toprule
Meilleur cas & $\Th(n\log n)$\\
Cas moyen & $\Th(n\log n)$\\
Pire cas & $\Th(n\log n)$ --- \textbf{aucune entrée ne le ralentit}\\
Mémoire auxiliaire & $\Th(n)$ : \textbf{n'est pas un tri en place}\\
Stabilité & \textbf{stable} grâce au test \tab{u[i] <= v[j]}\\
\bottomrule
\end{tabular}
\end{center}
\vspace{8pt}
\begin{encadre}\small
\textbf{Atout.} Une complexité garantie et la stabilité : c'est le tri de référence
lorsqu'on veut une performance sûre, ou lorsque les données sont lues séquentiellement
(fichiers, listes chaînées).
\end{encadre}
\begin{attention}\small
\textbf{Limite.} Il faut recopier les données : $\Th(n)$ de mémoire supplémentaire.
\end{attention}
\end{frame}

% =====================================================================
\section{Le tri rapide}
% =====================================================================

\begin{frame}{Principe}
\begin{encadre}
On choisit un élément, le \motcle{pivot}. On \textbf{partitionne} le tableau :
à gauche les éléments plus petits que le pivot, à droite les plus grands.
Le pivot est alors \emph{à sa place définitive}, et il ne reste qu'à trier
récursivement les deux côtés --- \textbf{sans rien recombiner}.
\end{encadre}
\vspace{6pt}
\begin{center}
\begin{tikzpicture}[font=\scriptsize,>=Stealth]
  \draw[draw=teal,fill=grisclair] (0,0) rectangle (9,0.7);
  \node at (4.5,0.35) {tableau quelconque};
  \draw[->,thick,teal] (4.5,-0.1)--(4.5,-0.7) node[midway,right] {partition : coût $\Th(n)$};
  \draw[draw=teal,fill=tealclair] (0,-1.45) rectangle (3.6,-0.75);
  \node at (1.8,-1.1) {éléments $\le p$};
  \draw[draw=orangefonce,fill=orangeclair] (3.6,-1.45) rectangle (4.8,-0.75);
  \node at (4.2,-1.1) {\textbf{p}};
  \draw[draw=teal,fill=tealclair] (4.8,-1.45) rectangle (9,-0.75);
  \node at (6.9,-1.1) {éléments $>p$};
  \draw[->,thick,vertrec] (1.8,-1.6)--(1.8,-2.1) node[midway,left] {tri récursif};
  \draw[->,thick,vertrec] (6.9,-1.6)--(6.9,-2.1) node[midway,right] {tri récursif};
  \node[orangefonce] at (4.2,-2.2) {\textbf{définitivement placé}};
\end{tikzpicture}
\end{center}
\end{frame}

\begin{frame}[fragile]{Tri rapide --- implémentation}
\begin{lstlisting}
def tri_rapide(L):
    if len(L) <= 1:
        return L
    pivot = L[0]
    petits = [x for x in L[1:] if x <= pivot]
    grands = [x for x in L[1:] if x > pivot]
    return tri_rapide(petits) + [pivot] + tri_rapide(grands)
\end{lstlisting}
\begin{encadre}\small
Le pivot choisi est le \textbf{premier élément}. Les deux listes en compréhension
réalisent la partition : chacune parcourt $L[1:]$, soit $n-1$ comparaisons au total
pour l'une et $n-1$ pour l'autre, donc un coût $\Th(n)$ par appel.
La concaténation finale place le pivot \emph{définitivement} entre les deux parties.
\end{encadre}
\end{frame}

\begin{frame}{Terminaison et correction}
\begin{terminaisonrec}
La longueur $|L|$ est un entier positif. Comme \tab{petits} et \tab{grands} sont
extraits de $L[1:]$, on a $|\text{petits}|\le n-1$ et $|\text{grands}|\le n-1$ :
la taille décroît strictement à chaque appel, et le cas de base $|L|\le1$ est atteint.
\end{terminaisonrec}
\vspace{4pt}
\begin{correction}[Correction --- récurrence forte sur la longueur]
\small
\emph{Base :} une liste de $0$ ou $1$ élément est triée.\\[2pt]
\emph{Hérédité :} par construction, tout élément de \tab{petits} est $\le$ \tab{pivot}
et tout élément de \tab{grands} est $>$ \tab{pivot}. Par hypothèse de récurrence,
\tab{tri\_rapide(petits)} et \tab{tri\_rapide(grands)} sont triées ; la concaténation
\tab{petits\_triés + [pivot] + grands\_triés} est donc triée, et elle contient
exactement les mêmes éléments que $L$ (chaque élément de $L[1:]$ va dans
\emph{exactement} une des deux listes).
\end{correction}
\end{frame}

\begin{frame}[shrink=8]{Complexité : tout dépend du pivot}
\begin{encadre}[title=Le meilleur cas : le pivot coupe en deux,fonttitle=\bfseries]
$C(n)=2\,C(n/2)+n$ \; d'où \; $C(n)=\Th(n\log n)$, comme le tri par fusion.
\end{encadre}
\vspace{4pt}
\begin{attention}
\textbf{Le pire cas : le pivot est toujours le plus grand (ou le plus petit)}\\[3pt]
C'est ce qui arrive si la liste est \emph{déjà triée} : une liste \tab{petits} vide et une liste \tab{grands}
de taille $n-1$.
\[
C(n)=C(n-1)+(n-1)\quad\Longrightarrow\quad C(n)=\frac{n(n-1)}2=\Th(n^2).
\]
\end{attention}
\vspace{4pt}
\begin{encadre}[title=Le cas moyen,fonttitle=\bfseries]
Sur une entrée aléatoire, on démontre que
$C_{\text{moy}}(n)\approx 1{,}39\,n\log_2 n=\Th(n\log n)$ :
seulement $39\,\%$ de plus que l'optimum.
\end{encadre}
\end{frame}

\begin{frame}{Le choix du pivot}
Le pire cas vient d'un pivot mal choisi. Trois parades classiques :
\vspace{4pt}
\begin{encadre}
\textbf{1. Pivot aléatoire} --- remplacer \tab{pivot = L[0]} par
\tab{pivot = L[random.randrange(len(L))]}. Aucune entrée particulière n'est plus
« mauvaise » qu'une autre : le $\Th(n\log n)$ devient vrai
\emph{en moyenne pour toute entrée}.
\end{encadre}
\vspace{3pt}
\begin{encadre}
\textbf{2. Médiane de trois} --- prendre comme pivot la médiane du premier, du
dernier et de l'élément central de \tab{L}. Élimine le pire cas « liste déjà triée »,
de loin le plus fréquent en pratique.
\end{encadre}
\vspace{3pt}
\begin{encadre}
\textbf{3. Bascule} --- pour une liste de moins de $15$ éléments, terminer par un
tri par insertion, plus rapide sur les petites tailles.
\end{encadre}
\end{frame}

\begin{frame}{Tri rapide : bilan}
\begin{center}\small
\begin{tabular}{ll}
\toprule
Meilleur cas & $\Th(n\log n)$\\
Cas moyen & $\Th(n\log n)$, environ $1{,}39\,n\log_2 n$\\
Pire cas & $\Th(n^2)$ (tableau déjà trié, pivot extrême)\\
Mémoire auxiliaire & $\Th(n)$ : les listes \tab{petits} et \tab{grands} sont recopiées\\
Stabilité & \textbf{non stable}\\
\bottomrule
\end{tabular}
\end{center}
\vspace{8pt}
\begin{attention}\small
\textbf{À ne pas dire.} « Le tri rapide est en $\Th(n\log n)$ » est incomplet :
il l'est \textbf{en moyenne}, et il est en $\Th(n^2)$ dans le pire cas.
C'est pourtant, en moyenne, l'un des plus rapides en pratique.
\end{attention}
\end{frame}

% =====================================================================
\section{Comparaison}
% =====================================================================

\begin{frame}{Comparaison des cinq algorithmes}
\begin{center}\scriptsize
\begin{tabular}{lcccccc}
\toprule
\textbf{Tri} & \textbf{Meilleur} & \textbf{Moyen} & \textbf{Pire} & \textbf{Mémoire} & \textbf{En place} & \textbf{Stable}\\
\midrule
Sélection & $\Th(n^2)$ & $\Th(n^2)$ & $\Th(n^2)$ & $\Th(1)$ & oui & non\\
Insertion & $\Th(n)$ & $\Th(n^2)$ & $\Th(n^2)$ & $\Th(1)$ & oui & oui\\
Bulles & $\Th(n)$ & $\Th(n^2)$ & $\Th(n^2)$ & $\Th(1)$ & oui & oui\\
\midrule
\rowcolor{tealclair}
Fusion & $\Th(n\log n)$ & $\Th(n\log n)$ & $\Th(n\log n)$ & $\Th(n)$ & non & oui\\
\rowcolor{tealclair}
Rapide & $\Th(n\log n)$ & $\Th(n\log n)$ & $\Th(n^2)$ & $\Th(n)$ & non & non\\
\bottomrule
\end{tabular}
\end{center}
\vspace{8pt}
\begin{encadre}\small
\textbf{Ordres de grandeur pour $n=10^6$ :} environ $5\cdot10^{11}$ comparaisons pour un
tri quadratique (plus d'une heure) contre $2\cdot10^{7}$ pour un tri en $n\log n$
(une fraction de seconde).
\end{encadre}
\end{frame}

\begin{frame}{À retenir}
\begin{encadre}[title=Quatre points,fonttitle=\bfseries]
\begin{enumerate}\setlength\itemsep{3pt}
  \item Les deux algorithmes appliquent \motcle{diviser pour régner}, mais placent
        l'effort à des endroits opposés : la fusion pour l'un, la partition pour l'autre.
  \item Le \motcle{tri par fusion} est en $\Th(n\log n)$ \emph{toujours}, et il est
        \emph{stable} ; en revanche il n'est pas en place ($\Th(n)$ de mémoire).
  \item Le \motcle{tri rapide} est en $\Th(n\log n)$ \emph{en moyenne} et très
        rapide en pratique, mais il est en $\Th(n^2)$ dans le pire cas et il n'est
        pas stable.
  \item Le \motcle{choix du pivot} (aléatoire, médiane de trois) est ce qui rend le
        tri rapide fiable.
\end{enumerate}
\end{encadre}
\end{frame}

\begin{frame}{Pour s'entraîner}
\begin{enumerate}\small\setlength\itemsep{5pt}
  \item Dérouler à la main le tri par fusion, puis le tri rapide, sur le tableau
        $[5,\,3,\,8,\,1,\,9,\,2]$. Compter les comparaisons dans chaque cas.
  \item Écrire une version de \tab{tri\_rapide} avec pivot aléatoire, et vérifier
        expérimentalement qu'un tableau déjà trié n'est plus un cas défavorable.
  \item Montrer sur un exemple que le tri rapide n'est pas stable.
  \item Comparer les temps d'exécution des cinq tris pour
        $n=10^2,10^3,10^4$ (module \tab{time}), et tracer les courbes obtenues.
  \item Modifier le tri par fusion pour qu'il compte en plus le nombre d'inversions
        du tableau initial, sans changer sa complexité.
\end{enumerate}
\end{frame}

\begin{frame}[plain]
  \vfill
  \begin{center}
    {\color{teal}\hrule height 2pt}\vspace{10pt}
    {\LARGE\bfseries\color{teal} Fin}\\[8pt]
    {\normalsize Tri rapide et tri par fusion --- CPGE MP / PSI / TSI}\\[10pt]
    {\color{teal}\hrule height 1pt}
  \end{center}
  \vfill
\end{frame}

\end{document}
